% ch-extending.tex -- how to add physics to Yuzen, as a worked procedure
\chapter{Extending the model}
\label{ch:extending}

This chapter is for a reader who wants Yuzen to solve something it does not
solve yet. It is written as a procedure: decide what kind of extension you
have, then follow the numbered steps for that kind, then test it. It does not
tour the architecture --- \cref{ch:layout} says where everything lives --- and
it does not re-derive the numerical method, which is
\cref{ch:discretization,ch:timeintegration}.

Two chapters are prerequisites rather than further reading.
\Cref{ch:reproducibility} states the four things in the code that keep results
bit-identical across thread counts; three of them constrain how you add a term,
and one of them --- the registration order --- is a step of the procedure below.
\Cref{sec:time:form,sec:time:jacobian} give the form every equation is cast
into and the rule for how much Jacobian a term needs. Everything here builds on
those and cites them rather than repeating them.

\section{Before you write any code}
\label{sec:extending:decide}

Three questions decide how much work you are in for. Answer them in this
order, because they are in descending order of cost.

\begin{description}
  \item[Does the effect enter as a flux contribution?] That is, does it add a
    term to $\Qvec$ in \cref{eq:flux} --- something proportional to a mobility
    times a gradient, or to a mobility times a body force? Then it is a
    \emph{physics term}, it belongs in the matrix $A$, and the main path of
    this chapter (\cref{sec:extending:term}) applies. This is the most work of
    the three: a kernel implementing seven members and its face flux, a dimensionless group, an
    input block, a velocity field and a test.
  \item[Is it a source or a sink in the height equation?] Something that adds
    or removes liquid without transporting it --- a second evaporation
    mechanism, a deposition rate, a chemical loss. It joins evaporation and the
    supply term in the vector $B$, assembled at the end of
    \code{WellMixed\_HeightEquationBase::SetMatrixA()}, and that is a handful
    of lines rather than a class. But read the warning of
    \cref{sec:time:form} before you choose this route: \textbf{only what goes
    into $A$ receives the $\theta$ weighting}, so a term placed in $B$ and
    evaluated at the old state alone is effectively backward Euler, and the
    time accuracy of the scheme drops to first order for that term. Average it
    over the step as the evaporation is, or it is cheap to write and expensive
    to justify.
  \item[Is it a change to the concentration equation or to a material
    property?] A new concentration dependence for the viscosity, the surface
    tension or the diffusivity; a different evaporation law; a new shape
    function. \textbf{This is the cheapest of the three by a wide margin}: it
    is one class or one branch, one or two input keys, and no matrix
    assembly, no Jacobian and no registration order to respect.
    \Cref{sec:extending:other} covers these.
\end{description}

If your effect is a pressure --- anything that can be written as
$\Qvec = M(H,C)\,\nabla p$ --- it is the first kind, and the disjoining
pressure is the worked example to copy. If it is a body force, uniform or not,
it is still the first kind, and in-plane gravity and the centrifugal term are
the examples.

\section{Adding a physics term, step by step}
\label{sec:extending:term}

\begin{enumerate}

\item \textbf{Write the dimensionless form and identify the group.}

Start on paper, not in the editor. Put the new contribution into the flux in
the same shape as the existing ones (\cref{ch:terms}), non-dimensionalize it
with the scales of \cref{sec:governing:scales}, and read off the single
dimensionless group that multiplies it. Every existing term has exactly one:
$\As$ for capillary pressure, $\Bo$ for vertical gravity, $\Spn$ for the
centrifugal term, $\Dpn$ for the disjoining pressure. If your term needs two
groups, one of them is probably a property ratio that belongs in
\file{FluidProperty.*} instead.

Decide the sign here as well. The governing equation is written
$\partial H/\partial T = -\nabla\cdot\Qvec$, so expanding the divergence fixes
the sign the coefficient must carry; vertical gravity ends up with
$-\Bo/3$ and the disjoining pressure with $+\Dpn/3$. That sign lives in
\code{\_Magnitude} and nowhere else, which is what lets the mixed velocity be a
plain sum (\cref{sec:verification:velocity}).

\item \textbf{Add the group to \code{NonDimensionalNumbers}.}

In \file{yuzen/src/model/NonDimensionalNumbers.h} add the member and its
accessor alongside \code{Bo()}, \code{Sp()} and \code{Dp()}; in
\file{NonDimensionalNumbers.cpp} compute it inside \code{Build()} from the
dimensional input, and add a line for it to \code{ToString()} so that it
appears with the others in the run log on standard output
(\cref{sec:output:running}).

\textbf{The group must be exactly zero when the term is disabled.} The existing
code does this in one of two ways, and either is fine. $\Dpn$ is set to
\val{0} first and computed only inside \code{if(inp.Physics.\bk
DisjoiningPressure.\bk Enable)}; $\Bo$, $\Pgx$, $\Pgy$, $\Spn$ and $\Ma$ are
computed unconditionally and then forced to zero at the end of \code{Build()}
when their \key{Enable} key is false.

Why this matters is worth stating plainly, because it is the whole basis of the
claim that adding a term changes nothing. A disabled term costs nothing: its
\code{SetMatrixA} is never called at all, because
\code{WellMixed\_HeightEquationBase} skips any term whose \code{Enabled()} is
false. And a disabled term can change nothing even if it were called, because
\code{\_Magnitude} is built from a group that is exactly \val{0}. Together with
the appending rule of step 5, that is what makes every case without your term
\textbf{bit-identical} before and after your change
(\cref{sec:repro:order}).

\item \textbf{Add the input block and its keys.}

Two files. In \file{yuzen/src/IO/Input.h}, add a nested class inside
\code{I\_Physics} with one member per key, mirroring
\code{I\_DisjoiningPressure}. In \file{yuzen/src/IO/Input.cpp}, add the block
to \code{ReadPhysics()} following the pattern the other blocks use:

\begin{itemize}
  \item ask for \code{getBlockSize} inside a \code{try}\,/\,\code{catch(const
    BlockNotFound\&)} and treat a missing block as \code{bnum = 0};
  \item call \code{Error} if \code{bnum > 1}, since a physics block may not
    repeat;
  \item read each key with \code{getin("Name", member, default)}, where
    \textbf{omitting the third argument makes the key required};
  \item give the \code{else} branch --- the one taken when the block is absent
    entirely --- an explicit value for every member, as the
    \key{DisjoiningPressure} block does. If you leave it out, the members keep
    whatever the constructor left there;
  \item check ranges with \code{CheckRange} and string-valued keys against a
    candidate list with \code{CheckCandidates}, calling \code{Error} on
    failure. A check that spans several blocks does not belong here; it belongs
    in \file{yuzen/src/IO/CheckParameters.cpp}.
\end{itemize}

\begin{warn}[The default in the \code{else} branch is what your users will get]
  A misspelled block or key stops the run (\cref{sec:inputformat:unknown}), but
  a block that is simply left out is not an error: a user who never writes
  your block gets the values of the \code{else} branch in silence, with no
  warning anywhere. Choose them so that the silent outcome is the term
  switched \emph{off}, as it is for every other \key{Physics} sub-block.
\end{warn}

Then document the keys in \cref{ch:ref-physics} in the same change, with the
type, the default and the units.

\item \textbf{Write the kernel.}

One header in \file{yuzen/src/model/well\_mixed/}, a template on the traits
class, deriving from \code{PhysicsTermBase<T>}
(\file{PhysicsTermBase.h}). Templating on the traits is what lets one body
serve both dimensions: everything that differs between 1-D and 2-D --- the
number of directions, the stencil, the spacing, the site-to-row map, the two
direction orderings --- is reached through \code{T}
(\file{GridTraits.h}), so the kernel never mentions a dimension.

If your term is a body force whose driving field has no gradient in it, do not
write the seven members at all: derive from \code{BodyForceKernel<T>} and
supply only \code{Drive}, \code{DriveFace}, \code{Weight}, \code{Mag},
\code{SideOutflow}, \code{VelocityField} and \code{VelocitySign}.
\file{CentrifugalKernel.h} is 155 lines including the license header, and the
only thing specific to the centrifugal force in it is that the driving field is
the distance from the rotation axis. If your term is a mobility times the
gradient of something, copy the skeleton of \file{VerticalGravityKernel.h}
instead, and implement the seven members of \cref{sec:extending:contract}
and its face flux (\code{FaceFlux}, after them; \code{BodyForceKernel}
provides it for a body force).

Four more places have to learn about the class. Add a \code{typedef} to
\textbf{both} \code{WM1D\_Traits} and \code{WM2D\_Traits} in
\file{WellMixedTraits.h}; add a member of that type to
\code{WellMixed\_HeightEquationBase} in \file{HeightEquationBase.h}; call its
\code{Build} in \code{WellMixed\_HeightEquationBase::Build}, next to the
others. And if the term produces a velocity --- it does, unless it is a pure
sink --- add a field for it to \code{Velocity\_WM1D} and
\code{Velocity\_WM2D} in \file{model/LubFields.h}, add it to the sum in
\code{Cal\_MixedVelocity()} in \file{LubFields.cpp}, and add a
\code{SetVelocityBC} line for it in
\code{WellMixed\_HeightEquationBase::CalVelocity()}.

\begin{warn}[The mixed velocity is a plain sum and must stay one]
  Add your field as one more \code{+} in \code{Cal\_MixedVelocity()}. Do not
  write the composition as a chain of differences: the signs are already
  correct inside each term's \code{\_Magnitude}, and a difference chain flips
  the sign of every term added after it. That is a defect this code has had,
  and \file{yuzen/tests/test\_mixed\_velocity.cpp} exists to stop it coming back
  (\cref{sec:verification:velocity}).
\end{warn}

\item \textbf{Register it at the end of \code{RegisterPhysics()}.}

\code{WellMixed\_HeightEquationBase::RegisterPhysics()} in
\file{HeightEquationBase.h} pushes the terms into \code{\_physics} in a fixed
order, and that order is the order of the \code{+=} operations into the matrix,
which is the column insertion order of the CRS rows, which fixes the rounding
of every sum taken over a row. \textbf{Append your
\code{\_physics.push\_back(\&\_YourTerm)} after the last line; do not insert
it, and do not tidy the list into alphabetical order.} Appending is what makes
the existing results bit-identical, and \cref{sec:repro:order} is the argument.

\item \textbf{Add the source file to \file{CMakeLists.txt}, if there is one.}

A kernel that is a header-only template needs no entry: none of the seven
kernel headers appears in \file{CMakeLists.txt}. If your term has a
dimension-specific \file{.cpp} --- as the capillary pressure term does, with
\file{LaplacePressure\_1D.cpp} and \file{LaplacePressure\_2D.cpp} --- add it to
\code{YUZEN\_CORE\_SRC} in the top-level \file{CMakeLists.txt}. A new test file
goes into \code{YUZEN\_TEST\_SRC} in \file{yuzen/tests/CMakeLists.txt}.

\item \textbf{Write the test.}

\Cref{sec:extending:testing} says what a good one looks like. Write it before
you run a case, because a term that is merely plausible on a picture is a term
that is wrong in a way you will find in six months.

\end{enumerate}

\section{What the seven members must do}
\label{sec:extending:contract}

The contract is declared in
\file{yuzen/src/model/\bk well\_mixed/\bk PhysicsTermBase.h}. Four members are
public and three are protected; all seven are pure virtual, so the compiler
will not let you forget one.

The distinction that runs through the whole contract is between
\textbf{coefficients that depend only on the mesh and the scheme}, which are
built once in \code{Build()} and held in the arrays \code{C(d)} and
\code{CV(a)}, and \textbf{quantities that depend on the field}, which are
recomputed from scratch on every call because the field changes at every outer
iteration. Putting a field-dependent quantity in the first group is a bug that
shows up as a Newton iteration that will not converge; putting a mesh-only
coefficient in the second group is merely slow.

\begin{description}

  \item[\code{Build(inp\_phys, g, nonDim, bc)}] Called once, from
    \code{WellMixed\_HeightEquationBase::Build}, before anything is solved. It
    is handed the \key{Physics} block of the input, the grid, the set of
    dimensionless numbers and the boundary condition. It must store those
    pointers, set \code{\_enabled} from your \key{Enable} key, set
    \code{\_Magnitude} to the coefficient of the whole term \emph{with its
    sign} --- the group over 3 for a term carrying the mobility $H^{3}/3$, and
    $\As\,\Ma/2$ for solutal capillarity, which carries $H^{2}/2$ --- call
    \code{T::SetExtent} to record the grid extent, and finish by calling
    \code{SetLinearCoef()}. It must not read the field: at this point the
    initial condition is not in it. This is also the only opportunity to read
    the input, so anything from the input file that the term needs must be
    copied into a member here.

  \item[\code{SetLinearCoef()}] Protected; called from \code{Build()} and from
    nowhere else. It fills \code{C(d)[s]}, the difference coefficient of
    stencil point \code{s} for direction \code{d}, and \code{CV(a)[s]}, its
    counterpart for the velocity along axis \code{a}. These depend on the mesh
    spacing \code{T::Spacing} and on the scheme, and on nothing else.
    \textbf{It must not look at the field}, which is what makes it safe to
    call once.

  \item[\code{SetMatrixA(crsA, field)}] Called on every outer iteration
    (\cref{sec:time:newton}), with \code{field} the current Newton guess, not
    the committed field. It must \textbf{return immediately if
    \code{!\_enabled}}, skip every row for which \code{crsA.IsBC(irow)} is
    true, and \textbf{add} its contribution with \code{+=} --- through
    \code{crsA.Diag(irow)} for the diagonal and \code{crsA(irow, col)}
    otherwise. It must never assign, never clear the matrix, and never write to
    $B$: the clearing is done for it, and $B$ is assembled afterwards by the
    height equation. The sum over directions within a row is taken in
    \textbf{axis order}, \code{T::Dir(d)} (\cref{sec:repro:directions}).

  \item[\code{SetJacobiMatrix(crsJ, field)}] Called once every
    \key{JacobianUpdateInterval} iterations (\cref{sec:time:jacobian}),
    after $A$ has already been copied into the Jacobian; your term adds its own
    derivative $\partial(A\varphi)/\partial\varphi$ on top. The same rules
    apply --- return when disabled, skip boundary rows, add rather than assign
    --- with one difference that matters: the sums that land on the
    \textbf{diagonal} are taken in \textbf{stencil order},
    \code{T::StencilDir(k)}, not in axis order. \textbf{This member is the one
    that is allowed to be approximate}; see \cref{sec:extending:jacobian}.

  \item[\code{NonlinearTerm(site, tedge, field)}] Protected. The
    field-dependent factor on a face --- for almost every term the mobility
    $H^{3}/\hat\mu(C)$, averaged across the face, with the viscosity evaluated
    at the mean concentration of the two cells. \code{tedge} selects the face;
    the value \code{T::DirCenter()} (\code{OO}) means the cell itself rather
    than a face. It is recomputed on every call and must be a pure function of
    the field: no caching, no state.

  \item[\code{NonlinearJacobi(site, tedge, adj, field)}] Protected. The
    derivative of \code{NonlinearTerm} with respect to the height at one
    stencil point, \code{adj} saying which point. It is the second member that
    may be approximate, and dropping a piece of it costs iterations and not
    accuracy.

  \item[\code{CalVelocity(field)}] Called from
    \code{WellMixed\_HeightEquationBase::CalVelocity()}, once at the start of
    the run and once per step, for the state the step ends on. It writes the velocity
    \emph{this term alone} produces into \emph{its own} field --- the flux
    through a face, divided by \code{T::FaceHeight} --- and leaves every site
    for which \code{T::InsideForVelocity} is false alone. \textbf{It must not
    touch the mixed velocity and must not accumulate into another term's
    field.} That is precisely why the order of the terms does not matter here,
    unlike in \code{SetMatrixA} (\cref{sec:repro:order}). The mixed velocity is
    formed afterwards, as the plain sum of the individual fields; it is what
    the output reports, and what the Courant check is formed from. (What
    crosses an open edge is counted from the face fluxes, \code{FaceFlux}
    below.)

\end{description}

Two more members carry the term's flux to the well-mixed concentration. They
are not pure virtual --- \code{PhysicsTermBase} returns 0 for both --- so
\textbf{the compiler will not remind you of them}, and a term that leaves them
out moves the film without moving the solute with it.

\begin{description}
  \item[\code{FaceFlux(site, d, field)}] The volume flux of the term through
    face \code{d} of the cell, in the state \code{field}, exactly as its
    entries of $A$ carry it: along the axis of the face, per unit cell measure
    (divided by the spacing across the face), so that the row of the cell takes
    $\sum_d (\pm)$\code{FaceFlux} --- minus for the low faces, plus for the
    high ones --- which is its $A h$ to rounding. It must return 0 when the term
    is disabled, and 0 through a wall face where $A$ has no flux. The
    film-height equation forms each face's flux once from these
    (\code{FaceVolumeFluxes}), the concentration equation carries the
    solute with exactly that flux (\cref{sec:disc:wm-conservative}), and
    \file{History.dat} counts what crosses an open edge from it
    (\code{EdgeFluxRate}, \cref{sec:output:history}).
  \item[\code{SideRate(site, field)}] A volume sink of the cell that is not a
    flux through a face --- $A h$ of a diagonal term, such as the sideways
    outflow of the 1-D centrifugal term. The solute leaves with it at the
    concentration of the cell. Most terms have none.
\end{description}

The face fluxes are checked against $A h$ in \file{test\_jacobian.cpp}:
every term and every face interpolation, 1-D and 2-D, must agree with the
product of its entries to rounding, and the two cells that share a face must
see the same flux. Add your term to that list as you add it to the Jacobian's.

\section{How much Jacobian your term needs}
\label{sec:extending:jacobian}

The governing principle is \cref{sec:time:jacobian} in one sentence: an
approximate Jacobian changes the \emph{rate} of convergence and not the fixed
point, because what fixes the converged answer is the residual of
\cref{eq:nk-update} and not the matrix multiplying the correction.

The practical rule is to get it exactly right anyway, because it can be
checked exactly. \file{yuzen/tests/test\_jacobian.cpp}
(\cref{sec:verification:jacobian}) builds each term on a film, a substrate and
a concentration that all vary, and compares $J v$ with a central difference of
$A(h)h$ for a rough direction $v$; the two agree to $10^{-10}$ for every term
of the model. Add your term to its list: a Jacobian that leaves something out
shows up as the size of what it leaves out. The three faults it found were in
the vertical gravity, the disjoining pressure and the 1-D centrifugal
outflow, and they cost the disjoining-pressure cases 30--55 \% of their
Newton iterations.

Two things are easy to get wrong. The Jacobian starts as a copy of $A$, so a
term whose $A$ is not linear in $h$ --- the disjoining pressure writes
$\hat\Pi(h)$ into the matrix as $(\hat\Pi/h)\,h$ --- must add the derivative
\emph{less} what $A$ already holds. And the derivative of a face mobility
with respect to the cell itself is taken at the cell, not at the neighbor.

If you do use an approximate Jacobian, it changes the rate of convergence and
not the answer: watch the \code{NK\_iter} column of \file{CalcInfo.dat}
(\cref{sec:output:calcinfo}); the existing cases run between 3 and 6
(\cref{sec:verification:coupling}).



\section{The rules you must not break}
\label{sec:extending:rules}

\begin{warn}[Four rules, and one acceptance criterion]
  Each of these is argued in \cref{ch:reproducibility}; they are collected here
  because they are all constraints on the change you are about to make.

  \begin{itemize}
    \item \textbf{Append to \code{RegisterPhysics()}; never insert.} The order
      of that list is the column insertion order of the matrix and therefore
      the order of the floating-point additions (\cref{sec:repro:order}).
    \item \textbf{Do not touch the reduction scheme of \code{Vector}.} The
      fixed-block partial sums are what make an inner product independent of
      the thread count (\cref{sec:repro:reduction}).
    \item \textbf{Do not sort a CRS row into ascending column order.} Entries
      are stored in insertion order with the diagonal first, and libraries that
      need sorted indices are handed the sorted mirror
      (\cref{sec:repro:crs}).
    \item \textbf{Do not unify the two direction orderings.} \code{T::Dir(d)}
      is axis order and \code{T::StencilDir(k)} is stencil order; they are used
      deliberately in different places, and in 1-D they coincide, so the
      damage appears only when somebody runs a 2-D case
      (\cref{sec:repro:directions}).
  \end{itemize}

  \textbf{The acceptance criterion for a change that is meant to add a term
  without altering existing results is bit-identity.} Run the shipped cases
  before and after, with the term disabled, and compare with
  \file{tools/compare.py}: every column must read \code{0.000e+00}. Not ``small
  differences'', not ``differences at the level of the linear-solver tolerance''.
  Zero. If it is not zero, you have inserted rather than appended, or changed
  something that was already being summed.
\end{warn}

\section{Testing a new term}
\label{sec:extending:testing}

The repository has two tests worth studying before you write yours, and they
illustrate two different kinds of claim.
\file{yuzen/tests/test\_disjoining.cpp} checks a physical term against
mathematics that exists independently of the code.
\file{yuzen/tests/test\_mixed\_velocity.cpp} checks a composition rule with
values chosen so that any error identifies itself. Between them they suggest
three things to aim for, in descending order of strength
(\cref{sec:verification:policy}).

\begin{description}

  \item[A state the term must leave alone.] The cheapest strong test there is.
    For the disjoining pressure it is a uniform precursor film $H = \Hstar$,
    where $\Pihat(\Hstar) = 0$ identically, so the term must produce no
    velocity whatsoever. Find the analogous state for your term --- a uniform
    film for anything driven by a gradient, zero rotation rate for anything
    driven by rotation --- and require the contribution to be below
    $10^{-12}$. Copy one detail from that test: it \textbf{dirties the velocity
    field first}, filling it with 1.0 before calling \code{CalVelocity}, so
    that writing a zero is distinguished from never writing at all. A term that
    silently does nothing passes the naive version of this test perfectly.

  \item[A state with a closed-form answer.] Stronger, and usually available for
    a new term in some limit even when it is not available in general. The
    disjoining-pressure calibration is checked this way, with no numerics at
    all: the equilibrium first integral gives the far-field slope in closed
    form, and the code must agree with it to $10^{-14}$ relative over 36
    parameter combinations (\cref{sec:verification:disjoining}). Sweep
    parameters rather than testing one set: an algebraic slip in a single
    factor survives one combination and does not survive a grid of them.

  \item[A refinement study.] Put an analytic state on the mesh, measure what
    the discretization leaves behind, and check that it falls at the order you
    intended when the mesh is refined. This is the test that catches a wrong
    coefficient, a wrong sign or a wrong stencil, all of which can survive a
    comparison on a single mesh. \Cref{tab:eqresidual} is the model: three
    meshes, the residual velocity falling by a factor of about four each time.
    Assert the ratio within a band rather than the order to two digits, as that
    test does --- a wrong term puts the ratio near 1, not near 3.5.

\end{description}

\begin{tip}[Add the term to the manufactured solution]
  \file{tools/mms\_generate.py} derives the manufactured sources from the
  governing equations with SymPy (\cref{sec:verification:mms}). Add the new
  flux to $\mathbf{Q}$ there, regenerate \file{ManufacturedSourceTerms.inc},
  add the parameters to \code{MmsParams} and to
  \code{C\_ManufacturedSource::Build}, and switch the term on in the
  \code{wm\_all} cases of \file{tools/mms.py}, tuned to move $h_t$ by about
  as much as the others. If it converges at second order there with
  everything else on, it is consistent with the equation it claims to
  discretize, in space and in time.
\end{tip}

\begin{warn}[Do not let a regression test be the only test]
  A test that runs the new term once and records what came out is
  \textbf{the weakest option of the four in \cref{sec:verification:policy}}. It
  says nothing at all about whether the behavior is right; it says only that
  it has not changed since the day you wrote it down --- including, if you
  wrote it down on a bad day, that it is reliably wrong. Regression tests are
  worth having, and the input-parsing and output-format tests in this
  repository are exactly that, but a new physical term whose only test is a
  regression test has not been verified in any useful sense.
\end{warn}

Finally, make sure the test does not depend on the thread count. Everything in
\file{yuzen/tests/} is run three times by \code{ctest} --- at the environment's
own thread count, at \env{OMP\_NUM\_THREADS=1} and at
\env{OMP\_NUM\_THREADS=4} --- so a test that compares against a value computed
by a reduction will fail in one of the three unless the invariance really
holds. That is the point of running it three times.

\section{Smaller extensions}
\label{sec:extending:other}

\subsection{A new shape function}

Shapes are used for both the initial free surface and the substrate profile.
Add the class to \file{yuzen/src/\bk commonFunctions/\bk ShapeFunctions.cpp}
and its header, following \code{C\_RectBump} or \code{C\_EllipseBump}: a
\code{Build} that stores the parameters together with the \key{Order} and the
create type, and an \code{operator()(px, py)} that returns the height at a
point. Then add the repeated input block in \file{IO/Input.cpp} and
\file{Input.h} --- shape blocks may appear any number of times, so they are
read with \code{getBlockSize} into a vector --- and add the loop that
constructs your class to \code{C\_GenerateShape::Build} in
\file{commonFunctions/GenerateShape.cpp}, pushing an entry onto
\code{\_shapeOrder} so that the shape takes part in the \key{Order} sorting.
Nothing else changes: the shape machinery does not touch the matrix.

\subsection{A new boundary condition}

Boundary conditions live in \file{yuzen/src/model/\bk BoundaryCondition.h} and
\file{.cpp}. The height equation is fourth order, so each boundary needs two
conditions, and they are imposed by writing the condition itself into the two
ghost rows of the matrix rather than by modifying the interior equations
(\cref{sec:disc:ghost}). Add your type to the \code{tBoundaryCondition} enum,
write the rows for each of the three fields \code{Height},
\code{HeightJacobi} and \code{Conc}, and add the string to the candidate list
in \file{IO/Input.cpp}.

\begin{warn}[Two of the existing conditions are verified only for structure]
  \val{"FreshInlet"} and \val{"Outlet"} have been checked for the structure of
  the matrix they produce --- indices in range, one equation per ghost row ---
  and \textbf{not} against any problem whose inflow or outflow is known
  analytically (\cref{sec:notverified:numerics}). Do not treat them as
  worked examples of a correct inflow condition. If you need one that is right,
  you will have to verify it yourself, and \cref{ch:notverified} is where the
  result of that work belongs.
\end{warn}

\subsection{A new concentration dependence for a material property}

\file{yuzen/src/model/FluidProperty.h} and \file{.cpp} hold the three forms
that exist: \code{ConstantProperty}, \code{ConcLinearProperty} (used for the
surface tension) and \code{ConcExponentProperty} (used for the viscosity and
the diffusivity). Each is built from a reference concentration, a reference
value and one parameter, and each returns 1 at the reference concentration ---
these are ratios, not \si{\pascal\second}. Add a class in the same shape, hold
it as a member of \code{NonDimensionalNumbers}, build it in that class's
\code{Build()}, and add the input key that selects which form is used. Nothing
downstream needs to change: the physics terms ask for
\code{\_NonDim->Viscosity(c)} and do not know which form answered.

\subsection{A new linear-solver backend}

Implement \code{ILinearSolverBackend} --- \code{Name}, \code{Build} and
\code{Solve} --- in \file{yuzen/src/\bk numerics/\bk LinearSolverBackend.cpp},
inside the anonymous namespace where \code{EigenBackend} and
\code{AmgclBackend} already live, and add a branch to
\code{MakeLinearSolverBackend()}. \textbf{Keep the new library's headers and
its \code{\#ifdef} inside that \file{.cpp}}, exactly as the AMGCL backend does:
that is why \file{LinearSolverBackend.h} says nothing about whether AMGCL was
built, and why a build without it needs no conditional compilation anywhere
else. Respect the contract of \cref{sec:linsolve:contract} --- in particular,
the solution vector is also the initial guess, and the matrix and right-hand
side must be restored before returning.

Then answer one question in the documentation, because users will ask it:
\textbf{is your backend deterministic under a change of thread count?} Run the
same problem at one, two and four threads and compare the solution, the
iteration count and the residual with \code{==}, as
\file{yuzen/tests/test\_linear\_solver.cpp} does. \val{"eigen"} is bit-identical;
\val{"amgcl"} is not, because its internal reductions are
non-deterministic (\cref{sec:linsolve:choosing}). Either answer is acceptable
for a backend --- the answer being unknown is not.

\section{Contributing a change back}
\label{sec:extending:submitting}

Keep this short and mechanical. Before you offer the change:

\begin{enumerate}
  \item Build and run the test suite, \code{ctest --test-dir build}. All three
    thread-count entries must pass.
  \item Run the shipped cases (\cref{ch:cases}), before and after your change.
  \item Compare the two with \file{tools/compare.py}, fixing the backend and
    the thread count across the comparison so that you are measuring your
    change and not the rounding of a different linear solver.
  \item \textbf{State in the change description whether the numbers moved, and
    why.} ``No case changed; \file{compare.py} reads \code{0.000e+00}
    throughout'' is the expected statement for a new term or a refactoring. If
    something did move, say which case, by how much, and which of the three
    situations of \cref{sec:policy:results} it is.
  \item Update this manual in the same change, and the shorter companion
    documents \file{docs/theory.md}, \file{docs/numerics.md} and
    \file{docs/validation.md} where they touch the same point. \textbf{The
    manual is the authoritative description}; a change that lands in the code
    and not in the manual leaves the two out of step. If the change alters a key or a default, update
    \cref{ch:ref-physics} or \cref{ch:ref-numerics} as well; if it alters a
    number, update \cref{ch:verification}, and if it adds something that is
    \emph{not} verified, say so in \cref{ch:notverified} rather than leaving
    the omission to be inferred.
  \item Add an entry to \cref{ch:revisions}, in the form
    \cref{sec:rev:keeping} specifies.
\end{enumerate}
